\section{Proofs for Section~\ref{sec:upper}}
\label{app:upper}

This appendix proves the results of Section~\ref{sec:upper}. We keep the
conventions of Section~\ref{sec:upper-conventions}. All norms of vectors are
Euclidean, and matrix and multilinear norms are the induced operator norms.
$\bar B(c,\varrho)$ denotes the closed Euclidean ball with center $c$ and
radius $\varrho$. For a polynomial $F$, $\norm F_1$ is the sum of the
absolute values of its coefficients.

\subsection{Derivative bounds}
\label{app:upper-bounds}

\begin{lemma}
\label{lem:upper-derivatives}
Let $L'\ge2$ be an integer and put $\Lambda=2^{10L'^2}$. Let
$F\in\Q[X_1,\ldots,X_n]$ have $n\le L'$, total degree at most $L'$, at most
$L'$ terms, and coefficients of absolute value at most $2^{L'}$. Then for all
$x,y\in\bar B(0,2^{5L'})$,
\[
 |F(x)|\le\Lambda,\qquad \norm{\nabla F(x)}\le\Lambda,\qquad
 \norm{\nabla^2F(x)}\le\Lambda,\qquad
 \norm{\nabla^2F(x)-\nabla^2F(y)}\le\Lambda\norm{x-y}.
\]
If $T=(T_1,\ldots,T_n)$ is a polynomial map whose components satisfy the same
assumptions, with at most $L'$ terms in total, then for all
$x,y\in\bar B(0,2^{5L'})$,
\[
 \norm{T(x)}\le\Lambda,\qquad \norm{J_T(x)}\le\Lambda,\qquad
 \norm{J_T(x)-J_T(y)}\le\Lambda\norm{x-y}.
\]
\end{lemma}

\begin{proof}
Put $\varrho_1=2^{5L'}$, so $|x_i|\le\varrho_1$ for every coordinate of a
point of the ball. Let $X^\alpha$ be a monomial with $|\alpha|\le L'$ and let
$\beta$ be a multi-index with $j=|\beta|\le3$. The derivative
$\partial^\beta X^\alpha$ is a monomial of degree at most $L'$ multiplied by a
product of at most $j$ factors from $\{1,\ldots,L'\}$, so
$|\partial^\beta X^\alpha(x)|\le L'^j\varrho_1^{L'}=L'^j2^{5L'^2}$ on the
ball. Since $\norm F_1\le L'2^{L'}\le2^{2L'}$, every partial derivative of
$F$ of order $j\le3$ is at most $2^{2L'}L'^j2^{5L'^2}$ there. A symmetric
$j$-linear form on $\R^n$ whose coefficients are bounded by $c$ has operator
norm at most its Frobenius norm, hence at most $n^{j/2}c$. With $n\le L'$
and $\log_2L'\le L'$, the derivative of order $j\le3$ has norm at most
\[
 2^{5L'^2+2L'+(9/2)\log_2L'}\le2^{5L'^2+(13/2)L'}\le2^{10L'^2},
\]
because $(13/2)L'\le5L'^2$ for $L'\ge2$. The Lipschitz bound for $\nabla^2F$
follows by integrating the third derivative along the segment from $y$ to
$x$, which lies in the ball.

For the map, let $C_i=\norm{T_i}_1$, so $\sum_iC_i\le2^{2L'}$. On the ball,
$\norm{T(x)}\le\sum_i|T_i(x)|\le2^{2L'}2^{5L'^2}$. Next,
$\norm{J_T(x)}\le\sum_{i,j}|\partial_jT_i(x)|\le n\sum_iC_iL'2^{5L'^2}
\le2^{5L'^2+2L'+2\log_2L'}$. Finally,
$\norm{\nabla\partial_jT_i(z)}\le\sqrt n\,C_iL'^22^{5L'^2}$ on the ball, so
\[
 \norm{J_T(x)-J_T(y)}\le\sum_{i,j}|\partial_jT_i(x)-\partial_jT_i(y)|
 \le n^{3/2}\Bigl(\sum_iC_i\Bigr)L'^22^{5L'^2}\norm{x-y}.
\]
All three constants are at most $\Lambda$ by the estimate above.
\end{proof}

\subsection{Algebraic separation}
\label{app:upper-separation}

We use the following special case of the one-block quantifier-elimination
theorem in the form stated by \citet[Theorem~2.27]{Basu2014}: one block of
$s$ existentially quantified variables, one free variable, and integer
coefficients. In the notation of that theorem this is the case $\omega=1$,
$k_1=s$, $\ell=1$ and coefficient ring $\Z$.

\begin{quote}
\emph{Imported theorem.} Let $\mathcal P\subset\Z[y_1,\ldots,y_s,z]$ be a
finite set of polynomials of total degree at most $d$ whose coefficients
have bit size at most $\tau$, and let $\Psi(y,z)$ be a quantifier-free
Boolean formula whose atoms are sign conditions on elements of $\mathcal P$.
Then the formula $\exists y\,\Psi(y,z)$ is equivalent over $\R$ to a
formula
\[
 \Theta(z)=\bigvee_i\bigwedge_j\Bigl(\bigvee_l
 \operatorname{sign}(P_{ijl}(z))=\sigma_{ijl}\Bigr),
 \qquad \sigma_{ijl}\in\{-1,0,1\},
\]
with polynomials $P_{ijl}\in\Z[z]$ of degree at most $d^{O(s)}$ whose
coefficients have bit size at most $\tau d^{O(s)}$.
\end{quote}

The implied constants are absolute. Fix an integer $c_0\ge1$ such that, for
all $s\ge1$, $d\ge2$ and $\tau\ge1$, the degrees are at most $d^{c_0s}$ and
the coefficient bit sizes at most $\tau d^{c_0s}$. Every algorithm below uses
this one fixed integer; nothing depends on the instance. The constant can be
read off from the complexity analysis of the elimination algorithm cited in
\citet{Basu2014}, but its numerical value plays no role.

\begin{proof}[Proof of Lemma~\ref{lem:separation}]
Let $\Theta(z)$ be the formula given by the imported theorem for
$\exists y\,\Phi(y,z)$, and let $\mathcal Q$ be the finite set of
polynomials $P_{ijl}$ that occur in it. The set defined by $\Theta$ is
$\{\alpha\}$.

First, some nonzero $P\in\mathcal Q$ vanishes at $\alpha$. Otherwise there is
an open interval $U\ni\alpha$ on which every nonzero element of $\mathcal Q$
has constant nonzero sign; the zero polynomials have sign $0$ everywhere.
Then every atom of $\Theta$ has the same truth value at each $z\in U$ as at
$\alpha$, so $\Theta(z)$ holds on all of $U$. This contradicts the fact that
$\Theta$ defines the single point $\alpha$. In particular $\alpha$ is
algebraic.

Write $P(z)=\sum_{i=0}^ep_iz^i$ with $e\le d^{c_0s}$ and integer
coefficients satisfying $|p_i|<2^\gamma$, where $\gamma=\tau d^{c_0s}$.
Suppose $\alpha\ne0$, and let $m$ be the least index with $p_m\ne0$. The
polynomial $Q(z)=P(z)/z^m$ has $Q(\alpha)=0$ and nonzero integer constant
term $p_m$; it is not constant, because it vanishes at $\alpha$. If
$|\alpha|<1$, then
\[
 1\le|p_m|=\Bigl|\sum_{i=m+1}^ep_i\alpha^{i-m}\Bigr|
 \le|\alpha|\sum_{i=m+1}^e|p_i|\le|\alpha|\,e\,2^\gamma .
\]
Since $e\le d^{c_0s}\le\gamma\le2^\gamma$, this gives
$|\alpha|\ge2^{-2\gamma}$. If $|\alpha|\ge1$, the same bound also holds.
This proves~(b).

For~(a), let $s,d,\tau\le L^c$. Using $\log_2L\le L$,
\[
 \log_2\bigl(2\tau d^{c_0s}\bigr)
 \le1+c\log_2L+c_0L^c\cdot c\log_2L
 \le1+cL+cc_0L^{c+1}\le c(c_0+2)L^{c+1},
\]
so~(b) gives $|\alpha|\ge2^{-2^e}$ with $e=c(c_0+2)L^{c+1}$.
\end{proof}

\subsection{Division elimination and shifted comparisons}
\label{app:upper-clearing}

\begin{proof}[Proof of Lemma~\ref{lem:denominator-clearing}]
An integer $a$ with binary digits $a_k\cdots a_0$ is built from $0$ and $1$ by
Horner's rule: start with $0$, and for each digit replace the current node
$v$ by $v+v$ and then add $1$ if the digit is one. A negative integer is
$0-|a|$. This uses $O(k)$ nodes. A rational constant $a/b$ with $b>0$ is
represented by the pair $(a,b)$. Every other node $v$ of $C$, with operands
represented by $(N_a,Q_a)$ and $(N_b,Q_b)$, receives the pair
\[
 \begin{array}{ll}
 a+b: & (N_aQ_b+N_bQ_a,\ Q_aQ_b),\\
 a-b: & (N_aQ_b-N_bQ_a,\ Q_aQ_b),\\
 a\times b: & (N_aN_b,\ Q_aQ_b),\\
 a/b: & (N_aQ_bN_b,\ Q_aN_bN_b).
 \end{array}
\]
Each rule adds at most four nodes and refers to existing nodes, so sharing
is preserved and the size is $O(\sigma+b)$. The construction inspects no
values.

Suppose every division of $C$ is by a nonzero value. We show by induction
along a topological order that $Q_v>0$ and that the value of $v$ is
$N_v/Q_v$. Constants satisfy this. For sums, differences and products the
claim follows from $Q_a,Q_b>0$. For a division, the divisor has value
$N_b/Q_b\ne0$, so $N_b\ne0$ and $Q_aN_b^2>0$; moreover
$N_aQ_bN_b/(Q_aN_b^2)=(N_a/Q_a)(Q_b/N_b)$. Since $Q_v>0$, the value of $v$
and $N_v$ have the same sign.
\end{proof}

When the hypothesis fails, the integer circuit is still defined, but a
denominator can be zero; every application below verifies the hypothesis.

\begin{lemma}[Shifted comparisons]
\label{lem:upper-shifted}
Let $g>0$ and let $\alpha,\hat\alpha$ be real numbers such that $\alpha=0$ or
$|\alpha|\ge g$, and $|\hat\alpha-\alpha|\le g/8$. Then
\[
\begin{array}{lll}
 \alpha>0\iff\hat\alpha-g/2>0, &
 \alpha\ge0\iff\hat\alpha+g/2>0, &
 \alpha=0\iff g^2/4-\hat\alpha^2>0,\\[2pt]
 \alpha<0\iff-\hat\alpha-g/2>0, &
 \alpha\le0\iff-\hat\alpha+g/2>0, &
 \alpha\ne0\iff\hat\alpha^2-g^2/4>0.
\end{array}
\]
Each of the six expressions has absolute value at least $3g/8$ in the first
two columns and at least $15g^2/64$ in the last column.
\end{lemma}

\begin{proof}
If $\alpha>0$, then $\alpha\ge g$ and $\hat\alpha\ge7g/8$, so
$\hat\alpha-g/2\ge3g/8$. If $\alpha\le0$, then $\hat\alpha\le g/8$ and
$\hat\alpha-g/2\le-3g/8$. If $\alpha\ge0$, then $\hat\alpha\ge-g/8$ and
$\hat\alpha+g/2\ge3g/8$; if $\alpha<0$, then $\alpha\le-g$ and
$\hat\alpha+g/2\le-3g/8$. The rows for $<$ and $\le$ follow by replacing
$\alpha,\hat\alpha$ with $-\alpha,-\hat\alpha$. If $\alpha=0$, then
$|\hat\alpha|\le g/8$ and $g^2/4-\hat\alpha^2\ge15g^2/64$. If $\alpha\ne0$,
then $|\hat\alpha|\ge7g/8$ and $g^2/4-\hat\alpha^2\le-33g^2/64$. The last
expression is the negative of the equality expression.
\end{proof}

\subsection{Local Newton convergence and elimination without pivoting}
\label{app:upper-newton}

\begin{lemma}
\label{lem:upper-pivots}
Let $A\in\R^{m\times m}$ satisfy $v^{\mathsf T}Av>0$ for every $v\ne0$; $A$
need not be symmetric. Then every leading principal minor of $A$ is
positive. Consequently Gaussian elimination without row or column exchanges
applies to $A$, all of its pivots are positive, and a system $Ax=c$ is solved
with $O(m^3)$ arithmetic operations in which every division is by a pivot.
\end{lemma}

\begin{proof}
Let $A_k$ be the leading $k\times k$ submatrix. It satisfies
$v^{\mathsf T}A_kv>0$ for $v\in\R^k\setminus\{0\}$ (pad $v$ with zeros). For
$\theta\in[0,1]$ the matrix $(1-\theta)I+\theta A_k$ has the same property and
is therefore nonsingular: a null vector $v$ would give
$v^{\mathsf T}((1-\theta)I+\theta A_k)v=0$. Its determinant is a continuous
nonvanishing function of $\theta$ that equals one at $\theta=0$, so
$\det A_k>0$. Elimination without exchanges does not change leading principal
minors, and after $k-1$ steps the leading $k\times k$ block is upper triangular
with the first $k$ pivots on its diagonal. Hence the $k$th pivot equals
$\det A_k/\det A_{k-1}>0$, with $\det A_0=1$. Forward elimination and back
substitution divide only by pivots.
\end{proof}

\begin{lemma}[Local Newton convergence]
\label{lem:upper-local-newton}
Let $T\colon\R^m\to\R^m$ be continuously differentiable with Jacobian $J$,
and let $T(p)=0$. Let $0<\mu\le\Lambda$ and $\rho=\mu/(4\Lambda)$. Assume that
for all $x,y\in\bar B(p,\rho)$ and $v\in\R^m$,
\[
 v^{\mathsf T}J(x)v\ge\mu\norm v^2,\qquad
 \norm{J(x)-J(y)}\le\Lambda\norm{x-y}.
\]
Let $x_0\in\bar B(p,\rho)$ and $x_{t+1}=x_t-J(x_t)^{-1}T(x_t)$. Then every
$x_t$ lies in $\bar B(p,\rho)$, every $J(x_t)$ has positive leading principal
minors, and
\[
 \norm{x_t-p}\le\frac{2\mu}{\Lambda}\,2^{-3\cdot2^t}\qquad(t\ge0).
\]
\end{lemma}

\begin{proof}
For $x\in\bar B(p,\rho)$ and $v\in\R^m$, $\norm{J(x)v}\norm v\ge
v^{\mathsf T}J(x)v\ge\mu\norm v^2$, so $J(x)$ is invertible with
$\norm{J(x)^{-1}}\le1/\mu$, and Lemma~\ref{lem:upper-pivots} applies to it.
Suppose $x_t\in\bar B(p,\rho)$ and put $e_t=\norm{x_t-p}$. The segment from
$p$ to $x_t$ lies in the ball, and
$T(x_t)=\int_0^1J(p+\theta(x_t-p))(x_t-p)\,d\theta$. Hence
\[
 x_{t+1}-p=J(x_t)^{-1}\int_0^1\bigl(J(x_t)-J(p+\theta(x_t-p))\bigr)
 (x_t-p)\,d\theta,
\]
and since $\norm{J(x_t)-J(p+\theta(x_t-p))}\le\Lambda(1-\theta)e_t$,
\[
 e_{t+1}\le\frac1\mu\int_0^1\Lambda(1-\theta)e_t^2\,d\theta
 =\frac{\Lambda}{2\mu}e_t^2.
\]
Put $q_t=\Lambda e_t/(2\mu)$. Then $q_0\le\Lambda\rho/(2\mu)=1/8$ and
$q_{t+1}\le q_t^2$, so $q_t\le2^{-3\cdot2^t}$ and $e_{t+1}\le q_te_t\le
e_t/8$. Thus all iterates stay in $\bar B(p,\rho)$, and
$e_t=(2\mu/\Lambda)q_t$ gives the bound.
\end{proof}

\subsection{Starting points}
\label{app:upper-warm}

For gradients we use convex value optimization in the following form, which
is the content of Lemma~\ref{lem:convex-value} that we need: given an
explicit sparse convex polynomial $f$, an explicit rational polyhedron $P$
contained in a box with rational bounds, and a rational $\eta>0$, a
deterministic algorithm computes, in time polynomial in the encoding length
and $\log(1/\eta)$, a rational point $x\in P$ with $f(x)\le\min_Pf+\eta$.

For maps that are not gradients there is no objective. We use the
central-cut ellipsoid method, in the following form.

\begin{lemma}[Rounded-cut warm start]
\label{lem:upper-warm-start}
Let $n\ge1$, and let $R_0\ge1$ and $0<\delta_0<\rho_0$ be rational. Let
$p\in\R^n$ be unknown with $\bar B(p,\delta_0)\subseteq
Q_0=[-R_0,R_0]^n$. Suppose a deterministic procedure, given a rational
$x\in Q_0$, either \emph{accepts} $x$, in which case $\norm{x-p}\le\rho_0$,
or returns a nonzero rational vector $a$ such that $a^{\mathsf T}(y-x)\le0$
for every $y\in\bar B(p,\delta_0)$. Then a deterministic algorithm finds an
accepted point, provided that every returned vector has encoding length
polynomial in the query length, $n$, and the encoding lengths of
$R_0,\rho_0,\delta_0$. It calls the procedure only on points of encoding length
polynomial in $n$ and the encoding lengths of $R_0,\rho_0,\delta_0$, and it
performs polynomially many further bit operations in these quantities.
\end{lemma}

\begin{proof}
Extend the procedure to all rational points: if $x\notin Q_0$, choose $i$
with $|x_i|>R_0$ and return $a=\operatorname{sign}(x_i)e_i$. For
$y\in\bar B(p,\delta_0)\subseteq Q_0$ this gives
$a^{\mathsf T}(y-x)\le R_0-|x_i|<0$. In all cases divide a returned vector by
its largest absolute entry, which preserves the cut and gives
$\norm a_\infty=1$.

Let $n=1$. Start with $I_0=[-R_0,R_0]$ and query its midpoint $x$. If $x$ is
not accepted, the cut shows that $\bar B(p,\delta_0)$ lies on one side of
$x$; keep that half of the interval. Each interval contains
$\bar B(p,\delta_0)$, whose length is $2\delta_0$, and the $j$th interval has
length $2^{1-j}R_0$. Hence a query is accepted after at most
$\lceil\log_2(R_0/\delta_0)\rceil+1$ steps, and all midpoints are rational
numbers of polynomial length.

Let $n\ge2$. Put $K=\bar B(p,\rho_0)$ and $K_1=\bar B(p,\delta_0)\subseteq K$.
Since $\norm p\le\sqrt nR_0$, $K$ is contained in the ball of radius
$R=nR_0+\rho_0$ around $0$. The extended procedure is an oracle of the kind
allowed in \citet[Remark~3.2.33]{GroetschelLovaszSchrijver1988}: for every
rational query $y$ and every rational tolerance, it either asserts that $y$
lies in $K$, which is true for accepted points, or returns $c$ with
$\norm c_\infty=1$ and $c^{\mathsf T}x\le c^{\mathsf T}y$ for all $x\in K_1$.
Run the central-cut ellipsoid method of
\citet[Theorem~3.2.1]{GroetschelLovaszSchrijver1988} with this oracle,
initial radius $R$ and volume bound $\varepsilon=(\delta_0/n)^n$. By that
theorem, it returns the queried center when the oracle asserts membership.
The remark retains this acceptance alternative and allows cuts valid only on
$K_1$. Thus the method stops either at an accepted point or with an ellipsoid
of volume less than $\varepsilon$ that contains $K_1$. The second alternative is
impossible: the cube $p+[-\delta_0/n,\delta_0/n]^n$ lies in $K_1$, because
each of its points is within $\sqrt n\,\delta_0/n\le\delta_0$ of $p$, and its
volume $(2\delta_0/n)^n$ exceeds $\varepsilon$. The method is
oracle-polynomial in $n$, $\langle R\rangle$ and $\langle\varepsilon\rangle$,
and all its centers have encoding length polynomial in these quantities
\citep[Lemma~3.2.8]{GroetschelLovaszSchrijver1988}.
\end{proof}

The rounding and enlargement steps of the cited method are what keep the
centers short while preserving containment of $K_1$; rounding an exact
ellipsoid update arbitrarily would not have this property. The procedure need
not decide membership in the unknown ball $K_1$. It only has to accept
correctly and to return cuts that keep $K_1$.

\subsection{Proof of Lemma~\ref{lem:newton-circuit}}
\label{app:upper-newton-proof}

\paragraph{Existence, uniqueness and size of the zero.}
In case~(a), strong convexity makes $f$ coercive and strictly convex, so it
has a unique minimizer $p$, which is the unique zero of $T=\nabla f$; and
$\nabla f$ satisfies the monotonicity inequality of case~(b) with the same
$\mu$. In case~(b), let $\varrho>\norm{T(0)}/\mu$ and let
$\Pi$ be the Euclidean projection onto $\bar B(0,\varrho)$. The continuous map
$x\mapsto\Pi(x-T(x))$ sends this ball to itself and has a fixed point $x^*$ by
Brouwer's theorem. The projection inequality at $x^*$ reads
$T(x^*)^{\mathsf T}(y-x^*)\ge0$ for all $y\in\bar B(0,\varrho)$. If
$\norm{x^*}=\varrho$, the choice $y=0$ contradicts
\[
 T(x^*)^{\mathsf T}x^*
 =\ip{T(x^*)-T(0)}{x^*}+T(0)^{\mathsf T}x^*
 \ge\mu\varrho^2-\norm{T(0)}\varrho>0 .
\]
Hence $x^*$ is interior, and choosing $y=x^*\pm\theta u$ for small $\theta>0$
and every direction $u$ gives $T(x^*)=0$. If $T(p)=T(p')=0$, monotonicity
gives $0\ge\mu\norm{p-p'}^2$, so the zero $p$ is unique. Applying the
inequality to $p$ and $0$ gives $\mu\norm p^2\le-T(0)^{\mathsf T}p$, so
\[
 \norm p\le\norm{T(0)}/\mu\le2^{2L}\cdot2^L=2^{3L},
\]
because $T(0)$ is the vector of constant terms of $T$ (in case~(a), of
linear coefficients of $f$) and $\mu\ge2^{-L}$. Finally, dividing
$\ip{T(x+\theta v)-T(x)}{\theta v}\ge\mu\theta^2\norm v^2$ by $\theta^2$ and
letting $\theta\to0$ gives $v^{\mathsf T}J_T(x)v\ge\mu\norm v^2$ for all
$x,v$.

\paragraph{Constants.}
Put $\Lambda=2^{10L'^2}$ and $\rho=\mu/(4\Lambda)$; note that
$\mu\le2^L\le\Lambda$ and $\rho<1$. Since $\norm p+1\le2^{5L'}$, the ball
$\bar B(p,\rho)$ lies in $\bar B(0,2^{5L'})$. On that larger ball,
Lemma~\ref{lem:upper-derivatives} applies to $f$ (case~(a)) or $T$ (case~(b))
and to every observable $h$ of encoding length at most $L'$: $J_T$ is
$\Lambda$-Lipschitz, $\norm{J_T}\le\Lambda$, $\norm T\le\Lambda$, and
$\norm{\nabla h}\le\Lambda$. In case~(a), $J_T=\nabla^2f$ and its Lipschitz
bound comes from the third derivatives of $f$.

\paragraph{Starting point.}
In case~(a), apply Lemma~\ref{lem:convex-value} to $f$ on the box
$[-2^{3L},2^{3L}]^n$, which contains $p$, with
$\eta=\mu\rho^2/2=\mu^3/(32\Lambda^2)$. The output $x_0$ satisfies
$f(x_0)-f(p)\le\eta$, and Lemma~\ref{lem:models-strong-convexity} gives
$\frac\mu2\norm{x_0-p}^2\le f(x_0)-f(p)$, so $\norm{x_0-p}\le\rho$.

In case~(b), apply Lemma~\ref{lem:upper-warm-start} with $R_0=2^{4L}$,
$\rho_0=\rho$ and
\[
 \delta_0=\frac{\mu^3\rho^2}{2\Lambda^3}=\frac{\mu^5}{32\Lambda^5}<\rho .
\]
The procedure evaluates $T(x)$ exactly at the rational query $x\in Q_0$; it
accepts if $\norm{T(x)}^2\le(\mu\rho)^2$ and otherwise returns $a=T(x)$.
Evaluation takes polynomial time because the degree is at most $L$. Since
$\norm p_\infty\le2^{3L}$ and $\delta_0<1$, the ball $\bar B(p,\delta_0)$ lies
in $Q_0$. If $x$ is accepted, then
$\mu\norm{x-p}^2\le\ip{T(x)-T(p)}{x-p}\le\norm{T(x)}\norm{x-p}$ gives
$\norm{x-p}\le\norm{T(x)}/\mu\le\rho$. If $x$ is rejected, then $T(x)\ne0$.
Every point of $Q_0$ lies in $\bar B(0,\sqrt n2^{4L})\subseteq\bar B(0,2^{5L'})$,
so $\norm{T(x)}=\norm{T(x)-T(p)}\le\Lambda\norm{x-p}$ and
$\norm{T(x)}\le\Lambda$. Rejection gives $\norm{x-p}>\mu\rho/\Lambda$, and for
$y\in\bar B(p,\delta_0)$,
\[
 T(x)^{\mathsf T}(y-x)=-\ip{T(x)-T(p)}{x-p}+T(x)^{\mathsf T}(y-p)
 <-\mu\Bigl(\frac{\mu\rho}{\Lambda}\Bigr)^2+\Lambda\delta_0
 =-\Lambda\delta_0<0 .
\]
So every rejection returns a cut that keeps $\bar B(p,\delta_0)$, and the
lemma yields a rational $x_0$ with $\norm{x_0-p}\le\rho$ in time polynomial
in $L'$. In both cases $x_0$ is printed in binary and has polynomially many
bits.

\paragraph{Newton circuit.}
Lemma~\ref{lem:upper-local-newton} applies on $\bar B(p,\rho)$ with $m=n$ and
the constants $\mu,\Lambda$. Put $t_*=e(L')+1$ and $\hat x=x_{t_*}$. The
circuit takes the coordinates of $x_0$ as rational constants. One Newton step
evaluates $T$ and $J_T$ at the current nodes and solves
$J_T(x_t)d=T(x_t)$ by Gaussian elimination without exchanges
(Lemma~\ref{lem:upper-pivots}); then $x_{t+1}=x_t-d$. For the evaluation,
note that every entry of $T$ and $J_T$ (in case~(a), of $\nabla f$ and
$\nabla^2f$) is a rational combination of monomials $X^\beta$ obtained from
the at most $L$ monomials of the input by removing at most two variables.
There are at most $L(1+n+n^2)$ such monomials, each of degree at most $L$,
and each is computed from the coordinate nodes by at most $L$ multiplications.
Thus one step adds $O(L^4+n^3)$ nodes and rational constants of polynomial
bit length, and the whole circuit has $O(L'^2(L^4+n^3))$ nodes. All divisions
are by elimination pivots, which are positive, or occur inside rational
constants. Earlier steps are referenced, never copied, so the expanded
fractions are never formed.

\paragraph{Accuracy.}
By Lemma~\ref{lem:upper-local-newton},
$\norm{\hat x-p}\le(2\mu/\Lambda)2^{-3\cdot2^{t_*}}$, and all iterates lie in
$\bar B(p,\rho)$. For an observable $h$ of encoding length at most $L'$, the
segment from $p$ to $\hat x$ lies in this ball, where
$\norm{\nabla h}\le\Lambda$. Writing $e=e(L')$,
\[
 |h(\hat x)-h(p)|\le\Lambda\norm{\hat x-p}\le2\mu\,2^{-6\cdot2^e}
 \le2^{L+1-6\cdot2^e}\le2^{-2^e-3}=g_{L'}/8,
\]
where the last inequality uses $5\cdot2^e\ge L+4$.

\paragraph{Separation.}
Let $\delta_T$ be the product of the denominators of the coefficients of $f$
(case~(a)) or of each component $T_i$ (case~(b)), and $\delta_h$ the product
of the denominators of the coefficients of $h$. Consider the formula with
quantified variables $X\in\R^n$ and free variable $z$,
\[
 \delta_T\,T_1(X)=0,\ \ldots,\ \delta_T\,T_n(X)=0,\qquad
 \delta_hz-\delta_hh(X)=0,
\]
with $\delta_T$ chosen per component in case~(b). Its polynomials have integer
coefficients: in case~(a) they are $\delta_Tc_\alpha\alpha_i$ with
$|\delta_Tc_\alpha\alpha_i|\le L2^{2L}<2^{3L}$, in case~(b) they are at most
$2^{2L}$, and those of the last equation are at most $2^{2L'}$. Hence their
bit sizes are at most $\tau=4L'$, their degrees at most
$d=\max\{2,\deg T,\deg h\}\le L'$, and there are $s=n\le L'$ quantified
variables. Since $p$ is the only zero of $T$, the formula defines the single
value $z=h(p)$. Lemma~\ref{lem:separation}(b) gives $h(p)=0$ or
\[
 |h(p)|\ge2^{-8L'\cdot L'^{c_0L'}}\ge2^{-2^{e(L')}},
\]
because $\log_2(8L')+c_0L'\log_2L'\le3L'+c_0L'^2\le(c_0+3)L'^2$ for
$L'\ge2$. This proves Lemma~\ref{lem:newton-circuit}.\qed

\subsection{Proofs of Theorems~\ref{thm:exact-upper} and~\ref{thm:upper-monotone}}
\label{app:upper-convex-proof}

\begin{proof}[Proof of Theorem~\ref{thm:exact-upper}]
\emph{Case~(a).} On input $(f,h,\mu,{\bowtie})$ of length $L$, construct the
circuit $\hat x$ of Lemma~\ref{lem:newton-circuit}(a) with $L'=L$. Append the
nodes of $\hat\alpha=h(\hat x)$, computed monomial by monomial, and the node
$g=g_L$, obtained from the constant $1/2$ by $e(L)$ squarings. Append the
expression of Lemma~\ref{lem:upper-shifted} for the relation ${\bowtie}$.
Apply Lemma~\ref{lem:denominator-clearing} and output the integer circuit
whose output is the numerator of the final node.

On a promised input, Lemma~\ref{lem:newton-circuit} gives $h(p)=0$ or
$|h(p)|\ge g$, and $|\hat\alpha-h(p)|\le g/8$, so the final rational node is
positive exactly when $h(p)\bowtie0$ by Lemma~\ref{lem:upper-shifted}. All
divisions in the rational circuit are by nonzero values: elimination pivots,
the constant $2$ in $1/2$, and denominators of rational constants. Hence
Lemma~\ref{lem:denominator-clearing} shows that the output integer has the
same sign. The construction makes no oracle calls and uses polynomial time.
The warm start of Lemma~\ref{lem:convex-value} is run with a clock equal to
its polynomial time bound, so that the algorithm also halts in polynomial
time on inputs that violate the promise; its output is then irrelevant.

\emph{Case~(b).} First check the certificate. The list $w$ must consist of
distinct monomials $X^\beta v_i$ and contain $v_1,\ldots,v_n$, and $M$ must be
a symmetric rational $m\times m$ matrix. Expand
$v^{\mathsf T}\nabla^2f(X)v=\sum_{i,j}v_iv_j\,\partial_i\partial_jf(X)$, a sparse
polynomial with at most $n^2$ times as many terms as $f$, and
$w^{\mathsf T}Mw=\sum_{k,l}M_{kl}w_kw_l$, a list of $m^2$ monomials; collect
terms and compare coefficients. Test $M\succ0$ by symmetric Gaussian
elimination in exact rational arithmetic: $M\succ0$ exactly when all pivots
are positive, and the pivots are ratios of leading principal minors, so all
intermediate numbers have polynomial length. If a check fails, output the
circuit $0$. Otherwise put $\mu=\det M/(\tr M)^{m-1}$. By
Lemma~\ref{lem:models-det-trace}, $0<\mu\le\lambda_{\min}(M)$. For real $X$
and $v$,
\[
 v^{\mathsf T}\nabla^2f(X)v=w^{\mathsf T}Mw\ge\mu\norm{w(X,v)}^2
 \ge\mu\norm v^2,
\]
because every $v_i$ is a coordinate of $w(X,v)$. Thus $(f,h,\mu)$ satisfies
the promise of case~(a), and its encoding length is polynomial in $L$; apply
case~(a).
\end{proof}

\label{app:upper-monotone-proof}
\begin{proof}[Proof of Theorem~\ref{thm:upper-monotone}]
Case~(a) is the same as in Theorem~\ref{thm:exact-upper}, using
Lemma~\ref{lem:newton-circuit}(b); existence and uniqueness of $p$ are part
of that lemma. The ellipsoid procedure is run with a clock as well; if it
ends without an accepted point, which cannot happen on promised inputs, the
algorithm uses $x_0=0$. In case~(b), the certificate is checked as before, with
$v^{\mathsf T}J_T(X)v=\sum_{i,j}v_iv_j\,\partial_jT_i(X)$. A valid certificate
gives $v^{\mathsf T}J_T(x)v\ge\mu\norm v^2$ for all $x,v$, and then
\[
 \ip{T(x)-T(y)}{x-y}
 =\int_0^1(x-y)^{\mathsf T}J_T\bigl(y+\theta(x-y)\bigr)(x-y)\,d\theta
 \ge\mu\norm{x-y}^2 ,
\]
which is the promise of case~(a).
\end{proof}

The Newton step in the monotone case solves the nonsymmetric system
$J_T(x_t)d=T(x_t)$ directly. One can instead solve the symmetric positive
definite system $J_T^{\mathsf T}J_Td=J_T^{\mathsf T}T$ by symmetric
elimination; both have the same solution. Symmetric elimination must not be
applied to $J_T$ itself.

\subsection{Proof of Proposition~\ref{prop:upper-slack-gap}}
\label{app:upper-slack}

Let $L$ be the input length, $D=\max\{2,\deg f,\deg h\}\le L$, and write the
rows of $A$ as $a_i^{\mathsf T}$. All steps before the final circuit are
ordinary rational computations of polynomial bit complexity; every
algorithmic step is run with a clock, so that inputs violating the promises
also finish in polynomial time.

\emph{Step 1: feasibility and a containing box.} Decide whether
$P=\emptyset$ by linear programming; if so, output the circuit $0$. Otherwise
compute a rational $\bar x\in P$ of polynomial encoding length.
Lemma~\ref{lem:models-strong-convexity} with $K=P$ gives
$\norm{p_P-\bar x}\le2\norm{\nabla f(\bar x)}/\mu<R_P$ with
$R_P=1+2\norm{\nabla f(\bar x)}_1/\mu$. Let
$B_P=\{x:\norm{x-\bar x}_\infty\le R_P\}$ and
$B_P'=\{x:\norm{x-\bar x}_\infty\le R_P+1\}$; then $p_P$ lies in the interior
of $B_P$, and every point within distance one of $p_P$ lies in $B_P'$.

\emph{Step 2: constants.} Put $R_2=\norm{\bar x}_\infty+R_P+1$. As in the
proof of Lemma~\ref{lem:upper-derivatives}, the rationals
\[
 \Lambda_f=n^2D^3\norm f_1R_2^D,\qquad
 \Lambda_h=nD\norm h_1R_2^D
\]
bound $\norm{\nabla^3f}$ and $\norm{\nabla h}$ on $B_P'$; they have
polynomial encoding length because $D\le L$. Multiply each row of $(A,b)$ by
the product of its denominators to obtain an integer matrix $(\bar A,\bar b)$
with the same rows up to positive factors. Each integer row of $\bar A$ has
Euclidean norm at most $2^{\ell_i}$, where $\ell_i$ is the encoding length of
row $i$, so by Hadamard's inequality every minor of $\bar A$ has absolute
value at most $2^L$. Consequently, for every set $I$ of independent rows and
every invertible column block $\bar A_{I,B}$, Cramer's rule bounds every entry
of $\bar A_{I,B}^{-1}\bar A_{I,N}$ by $2^L$. Row scaling by positive factors
cancels in this product, so $A_{I,B}^{-1}A_{I,N}=\bar A_{I,B}^{-1}\bar A_{I,N}$,
and the chart matrix $Z$ of Step~5 below satisfies $\norm Z\le\norm Z_F\le\zeta:=2^{2L+1}$. Put
\[
 \Lambda_\varphi=\zeta^3\Lambda_f+\zeta\Lambda_h+\mu\zeta,\qquad
 \rho_\varphi=\frac{\mu}{4\Lambda_\varphi},\qquad
 \sigma=\max\{1,\max_i\norm{a_i}_1\},
\]
\[
 \varepsilon=\min\Bigl\{1,\ \frac{\mu\delta^2}{128\sigma^2},\
 \frac{\mu\rho_\varphi^2}{2}\Bigr\}.
\]
Then $\zeta\rho_\varphi\le1/4$, and all these numbers have polynomial
encoding length.

\emph{Step 3: an approximate minimizer.} Apply Lemma~\ref{lem:convex-value}
to $f$ on the polyhedron $P\cap B_P$, whose minimizer is $p_P$, with accuracy
$\varepsilon$. The output $\tilde x\in P$ satisfies
$f(\tilde x)\le f(p_P)+\varepsilon$. Lemma~\ref{lem:models-strong-convexity}
with $K=P$ gives $\frac\mu2\norm{\tilde x-p_P}^2\le f(\tilde x)-f(p_P)$, hence
\[
 \norm{\tilde x-p_P}\le\sqrt{2\varepsilon/\mu}
 \le\min\{\delta/(8\sigma),\ \rho_\varphi\}.
\]

\emph{Step 4: the active set.} For every row,
$|a_i^{\mathsf T}(\tilde x-p_P)|\le\norm{a_i}_1\norm{\tilde x-p_P}\le\delta/8$.
A row that is active at $p_P$ therefore has slack at most $\delta/8$ at
$\tilde x$, and by the promise an inactive row has slack at least $\delta$ at
$p_P$ and at least $7\delta/8$ at $\tilde x$. Hence the active set is exactly
$S=\{i:b_i-a_i^{\mathsf T}\tilde x<\delta/2\}$, computed by exact rational
comparisons.

\emph{Step 5: restriction to the active affine space.} Let
$E=\{x:A_Sx=b_S\}$. If $d\in\ker A_S$, then $p_P\pm\theta d\in P$ for all
small $\theta>0$, because the rows in $S$ stay tight and the other rows have
positive slack. First-order optimality gives
$\pm\theta\nabla f(p_P)^{\mathsf T}d\ge0$, so $\nabla f(p_P)\perp\ker A_S$. Hence
$p_P$ is a stationary point of the strictly convex restriction of $f$ to $E$,
and it is the unique minimizer of $f$ on $E$.

Choose a maximal independent set $I\subseteq S$ of rows by rational
elimination; let $r=|I|$. Every row in $S$ is a linear combination of rows in
$I$, and $p_P$ satisfies all of them, so $E=\{x:A_Ix=b_I\}$. If $r=n$, then
$p_P=A_I^{-1}b_I$ is rational; evaluate $h(p_P)$ exactly and output the
circuit $1$ or $0$ according to whether $h(p_P)\bowtie0$. If $r<n$, choose
columns $B$ with $A_{I,B}$ invertible and let $N$ be the other columns. The
chart $x_N=y$, $x_B=A_{I,B}^{-1}(b_I-A_{I,N}y)$ writes $E=\{\bar x_E+Zy\}$,
where the rows of $Z$ indexed by $N$ form an identity matrix. Hence
$Z^{\mathsf T}Z\succeq I$. The function $\varphi(y)=f(\bar x_E+Zy)$ satisfies
$\nabla^2\varphi(y)=Z^{\mathsf T}\nabla^2f(\bar x_E+Zy)Z\succeq\mu I$, and its
unique minimizer is $y^*=(p_P)_N$. The case $S=\emptyset$ is covered with
$r=0$, $Z=I$ and $\bar x_E=0$.

\emph{Step 6: Newton steps in the chart.} Apply
Lemma~\ref{lem:upper-local-newton} to $T_\varphi(y)=Z^{\mathsf T}\nabla
f(\bar x_E+Zy)$ with $J_\varphi(y)=Z^{\mathsf T}\nabla^2f(\bar x_E+Zy)Z$, the
moduli $\mu$ and $\Lambda_\varphi$, and $x_0=y_0:=\tilde x_N$. Then
$\norm{y_0-y^*}\le\norm{\tilde x-p_P}\le\rho_\varphi$. For $y$ in
$\bar B(y^*,\rho_\varphi)$, the point $\bar x_E+Zy$ is within
$\zeta\rho_\varphi\le1/4$ of $p_P$ and so lies in $B_P'$. Therefore
$\norm{J_\varphi(y)-J_\varphi(y')}\le\zeta^3\Lambda_f\norm{y-y'}$ on that ball,
and the observable $\eta(y)=h(\bar x_E+Zy)$ has $\norm{\nabla\eta}\le
\zeta\Lambda_h\le\Lambda_\varphi$ there. The chart is not expanded into
monomials: one step computes $x=\bar x_E+Zy_t$, the sparse derivatives
$\nabla f(x)$ and $\nabla^2f(x)$ as in Appendix~\ref{app:upper-newton-proof},
the products with $Z$ and $Z^{\mathsf T}$, and one elimination without
exchanges, whose pivots are positive by Lemma~\ref{lem:upper-pivots}.

\emph{Step 7: separation.} Consider the formula with quantified variables
$(x,\lambda)\in\R^{n+r}$ and free variable $z$,
\[
 \bar A_Ix-\bar b_I=0,\qquad
 \delta_f\nabla f(x)-\bar A_I^{\mathsf T}\lambda=0,\qquad
 \delta_hz-\delta_hh(x)=0,
\]
where $\delta_f,\delta_h$ clear the denominators of $f$ and $h$. A pair
$(x,\lambda)$ satisfies the first two groups exactly when $x\in E$ and
$\nabla f(x)$ lies in the row space of $A_I$, that is, when $x$ is the
stationary point $p_P$ of $f$ on $E$; a multiplier exists at $p_P$. So the
formula defines the single value $h(p_P)$. It has $s=n+r\le2L$ quantified
variables, degree at most $D\le L$, and coefficient bit sizes at most $4L$.
Lemma~\ref{lem:separation}(b) gives $h(p_P)=0$ or
$|h(p_P)|\ge2^{-8L\cdot L^{2c_0L}}\ge g:=2^{-2^{e_2(L)}}$ with
$e_2(L)=(2c_0+3)L^2$.

\emph{Step 8: the sign.} Run $t_*=e_2(L)+1$ Newton steps and put
$\hat\alpha=\eta(y_{t_*})$. By Lemma~\ref{lem:upper-local-newton},
$|\hat\alpha-h(p_P)|\le\Lambda_\varphi\norm{y_{t_*}-y^*}\le2\mu\,2^{-6\cdot
2^{e_2(L)}}\le g/8$. Lemmas~\ref{lem:upper-shifted}
and~\ref{lem:denominator-clearing} give the integer circuit, as in the proof of
Theorem~\ref{thm:exact-upper}. All divisions in the rational circuit are by
positive pivots or inside constants, and its size is polynomial in $L$. The
set $S$ computed in Step~4 is the active set. \qed

\subsection{Proof of Theorem~\ref{thm:posslp-closure}}
\label{app:upper-closure}

\paragraph{Magnitudes.}
Order the nodes of $C$ topologically. On a Boolean input all node values are
integers. Inputs and constants have absolute value at most $2=2^{2^0}$. If
the first $k$ nodes are bounded by $\beta_k\ge2$ in absolute value, a sum or
difference is bounded by $2\beta_k\le\beta_k^2$, a product by $\beta_k^2$, and
a threshold value by $1$. Hence every value is bounded by
\[
 B=2^{2^S}.
\]
This uses binary gates and counts every node of $C$, including inputs,
constants, and the Boolean control logic of part~(b).

\paragraph{Range compression.}
For $\sigma\ge1$ let $F_\sigma(x)=2\sigma x/(\sigma+x^2)$. If
$1\le|x|\le\sigma$, then $F_\sigma(x)$ has the sign of $x$ and
$1\le|F_\sigma(x)|\le\sqrt\sigma$. Indeed, $F_\sigma$ is odd with positive
denominator; for $t\in[1,\sigma]$,
\[
 2\sigma t-(\sigma+t^2)=(t-1)(\sigma-t)+(\sigma-1)t\ge0,
\]
which gives $F_\sigma(t)\ge1$, and $\sigma+t^2\ge2\sqrt\sigma\,t$ gives
$F_\sigma(t)\le\sqrt\sigma$. Let $\sigma_j=2^{2^j}$, so
$\sqrt{\sigma_j}=\sigma_{j-1}$, and these constants are shared by repeated
squaring from $2$. If $1\le|x|\le\sigma_J$, applying
$F_{\sigma_J},F_{\sigma_{J-1}},\ldots,F_{\sigma_1}$ yields a number of the same
sign with absolute value in $[1,2]$.

\paragraph{Sign refinement.}
Let $G(x)=2x/(1+x^2)$. For $t\in[1,2]$, $G(t)\in[4/5,1]$. For $t\in[0,1]$,
$G(t)\ge t$ and
\[
 1-G(t)=\frac{(1-t)^2}{1+t^2}\le(1-t)^2 .
\]
So if $|x|\in[1,2]$ and $G$ is applied $k\ge1$ times, the result $y$ has
the sign of $x$ and $1-|y|\le(1/5)^{2^{k-1}}\le2^{-2^{k-1}}$.

\paragraph{Threshold replacement.}
Put $\varepsilon=2^{-2^{2S+4}}$. Let $v$ be the exact integer input of a
threshold gate and $\tilde v$ a rational number with $|\tilde v-v|\le1/4$.
Form $x=4\tilde v-2$. If $v\ge1$ then $x\ge1$, and if $v\le0$ then
$x\le-1$; also $|x|\le4B+3\le B^4=\sigma_{S+2}$, using $B\ge4$. Apply
$F_{\sigma_{S+2}},\ldots,F_{\sigma_1}$, then $G$ exactly $2S+5$ times, and
return $\tilde H=(1+y)/2$. By the two previous paragraphs,
$|\tilde H-H(v)|\le\frac12\,2^{-2^{2S+4}}\le\varepsilon$.

\paragraph{Error induction.}
Replace every arithmetic gate of $C$ by the same operation on the
approximate values, and every threshold gate by the replacement above,
applied to the approximate value of its input. Inputs and constants are
exact. Let $\Delta_k$ be the largest error among the first $k$ nodes and put
$\kappa=3B$. We show $\Delta_k\le\varepsilon\kappa^k$; it holds for input and
constant nodes. Note first that
\[
 \varepsilon\kappa^S<\tfrac14,
\]
because $\log_2(\varepsilon\kappa^S)\le-2^{2S+4}+S(2^S+2)<-2$ for $S\ge1$.
At an addition or subtraction the error is at most $2\Delta_{k-1}$. At a
multiplication with exact inputs $a,b$ and errors at most
$\Delta_{k-1}\le1$,
\[
 |\tilde a\tilde b-ab|\le|a|\,|\tilde b-b|+|\tilde b|\,|\tilde a-a|
 \le(2B+1)\Delta_{k-1}\le\kappa\Delta_{k-1}.
\]
At a threshold gate the input error is at most $\Delta_{k-1}<1/4$, so the
output error is at most $\varepsilon$. In all cases
$\Delta_k\le\varepsilon\kappa^k$. Errors are absolute, so cancellation and
reuse of a node twice are covered. The approximate output therefore lies
below $1/4$ when $C(u)=0$ and above $3/4$ when $C(u)=1$.

\paragraph{Positive denominators.}
Represent every approximate value by integer circuits $(P,Q)$ with $Q>0$,
starting from $(u_i,1)$ and the constants $(0,1),(1,1)$. Arithmetic gates use
the rules of Appendix~\ref{app:upper-clearing}, without division. The
threshold replacement uses
\begin{gather*}
 4\tilde v-2:\ (4P-2Q,\ Q),\qquad
 \frac{1+y}2:\ (Q+P,\ 2Q),\\
 F_\sigma(P/Q)=\frac{2\sigma PQ}{\sigma Q^2+P^2},\qquad
 G(P/Q)=\frac{2PQ}{Q^2+P^2},
\end{gather*}
and every new denominator is positive because $Q>0$ and $\sigma>0$. If the
final pair is $(P,Q)$, output the integer circuit $C'=2P-Q$. Then $C'(u)>0$
exactly when $P/Q>1/2$, that is, when $C(u)=1$. Each threshold gate costs
$O(S)$ nodes, the constants $\sigma_j$ cost $O(S)$ shared nodes, and each
arithmetic gate costs $O(1)$ nodes. So $C'$ has $O(S^2)$ nodes and is built in
time polynomial in $S$, without evaluating any value. This proves~(a).

\paragraph{Part (b): unrolling an oracle machine.}
Fix an encoding of integer circuits as binary strings, for instance a list of
instructions, each of which is a constant $0$ or $1$ or an operation
$(\circ,i,j)$ with $\circ\in\{+,-,\times\}$ and earlier indices $i,j$, the last
instruction being the output; a string that encodes no circuit is not in
$\PosSLP$. Let the machine use a query state in which it moves to a yes or no
state according to whether its oracle tape holds a string of $\PosSLP$. Let
$t=t(L)$ be a polynomial bound on its running time on inputs of length $L$,
valid for every sequence of oracle answers; if only a bound for the true
oracle is known, add a clock that rejects after $t$ steps, which does not
change the decided problem.

Encode the configuration after $\theta\le t$ steps by Boolean values: for each
tape, cell position at most $t$ and symbol, whether the cell holds the symbol;
for each tape and position, whether the head is there; and for each state,
whether it is the current state. The symbol under each head is the
disjunction over positions of ``the head is here and the cell holds this
symbol'', which takes $O(t)$ gates per tape. Every other value of the next
configuration is a Boolean function of the current state, the symbols under
the heads, a bounded number of neighboring cell and head values, and, in the
query state, one answer bit. Halting states are kept unchanged. Implement
Boolean operations by $\neg a=1-a$, $a\wedge b=ab$ and $a\vee b=a+b-ab$, which
are exact on bits. This gives $O(t)$ nodes per tape and step, plus the query
modules below. The input nodes are the $L$ input bits.

At every step $\theta$ attach a query module that computes the bit
$a_\theta=[\text{oracle tape content}\in\PosSLP]$; the transition uses
$a_\theta$ only when the state is the query state. The module works as
follows.
\begin{enumerate}[label=(\roman*)]
\item A string of length at most $t$ encodes at most $t$ instructions. A
deterministic polynomial-time algorithm maps the tape content to a validity
bit, and for each slot $k\le t$ to an opcode in
$\{\text{absent},0,1,+,-,\times\}$, two operand indices below $k$ in binary,
and the index of the output instruction. Unrolling this algorithm in the same
way as the machine gives a Boolean circuit of size polynomial in $t$.
\item Equality tests on the binary indices give bits $\epsilon_{k,i}$ and
$\epsilon'_{k,i}$ ($i<k$) for the first and second operand of slot $k$, bits
$o_k$ for the output slot, and one-hot opcode bits
$w_k$ (constant one), $s^+_k,s^-_k,s^\times_k$. All of them are $0$ for absent
or invalid slots.
\item Define integer nodes
\[
 Y_k=\sum_{i<k}\epsilon_{k,i}V_i,\quad Z_k=\sum_{i<k}\epsilon'_{k,i}V_i,\quad
 V_k=w_k+s^+_k(Y_k+Z_k)+s^-_k(Y_k-Z_k)+s^\times_kY_kZ_k ,
\]
and $a_\theta=\text{valid}\cdot H\bigl(\sum_ko_kV_k\bigr)$.
\end{enumerate}
If the string encodes a circuit, induction on $k$ shows that $V_k$ is the
value of instruction $k$, so $a_\theta$ is the correct answer. If it encodes
none, $a_\theta=0$, which is also correct. Every node value is an integer, and
the module has $O(t^2)$ nodes besides the parser. The module enumerates
addresses, not sequences of earlier answers, so the total number of nodes is
polynomial in $t$ although the queries depend on earlier answers. Products
$Y_kZ_k$ are formed even when they are not selected; they are nodes of the
circuit and are covered by the bound $B$.

The output of the unrolled circuit is the acceptance bit after $t$ steps, a
Boolean value. Part~(a) gives an integer circuit $C'$ of polynomial size with
$C'(u)>0$ exactly when the machine accepts $u$. Replacing the input nodes by
the constants given by the actual input yields one integer circuit, which is
the required $\PosSLP$ instance. On inputs satisfying the promise, the
machine answers correctly, and so does the instance. Conversely, a many-one
reduction to $\PosSLP$ is a $\mathrm P^{\PosSLP}$ algorithm with one query.

\paragraph{Part (c).}
Build one circuit containing $C_1,\ldots,C_k$ with their shared nodes, a
threshold gate on each output, and the gates of $\beta$ implemented
arithmetically as above. Its values are integers and its output is a bit, so
part~(a) applies. \qed
